%% kapa: exam=1 round=35 year=2024 period=2 form=A subject=appraisal_law subject_ko=감정평가관계법규 seq=39 number=39
%% source: Q-Net 감정평가사 기출문제 (cst00309) artlSeq=5213229 file=raw/1008/5213229/2교시 문제지 원본.pdf
%% conversion: kordoc 4.24.51 -> tools/gen1.py (auto); review=proofread-sonnet; flags=none
\begin{question}{39}
도시 및 주거환경정비법령상 정비구역에 관한 설명으로 옳은 것은?
\begin{choices}
  \item 광역시의 군수가 정비계획을 입안한 경우에는 직접 정비구역을 지정할 수 있다.
  \item 정비구역에서 건축물의 용도만을 변경하는 경우에는 따로 시장ㆍ군수등의 허가를 받지 않아도 된다.
  \item 재개발사업을 시행하는 지정개발자가 사업시행자 지정일부터 3년이 되는 날까지 사업시행계획인가를 신청하지 않은 경우 해당 정비구역을 해제하여야 한다.
  \item 토지등소유자는 공공재개발사업을 추진하려는 경우 정비계획의 입안권자에게 정비계획의 입안을 제안할 수 있다.
  \item 정비구역이 해제된 경우에도 정비계획으로 변경된 용도지역, 정비기반시설 등은 정비구역 지정 이후의 상태로 존속한다.
\end{choices}
\end{question}
\answer{4}
